%%%%%
%
%% Dans ce fichier, les �ventuels accents sont cod�s
%% en ISO-8859-15 (aka latin9)
%% La ligne suivante donne trois exemples 
%% e accent aigu : �, e accent grave : �, e-dans-l'o : �
\section{FSSP: Calcul de l'activit� � partir des comptages}

\itemstyle{Op\'erateur~:} \operfmt{fssp\_activ\_calc}
\index{fssp\_activ\_calc}

%%%%%OBJET%%%%
\itemstyle{Objet~:}
Calcul de la fraction d'activit� d'une sonde de type FSSP, � partir des comptages
dans la DOF ("\texttt{strobe}") et hors de la DOF ("\texttt{fastreset}").


%%%%%ARGUMENTS%%%%
\bigskip
%% Dans ce fichier, les �ventuels accents sont cod�s
%% en ISO-8859-15 (aka latin9)
%% La ligne suivante donne trois exemples 
%% e accent aigu : �, e accent grave : �, e-dans-l'o : �
%%%%%ENTREES%%%%
\begin{tabarg}
\multicolumn{3}{l}{\titretab {Entr\'ees}}\\
\hline
$Strobe$ & Vecteur, \texttt{float} & Nombre de particules dans la DOF 
(\textit{comptes par p�riode d'�chantillonnage})\\
$Fast\_Reset$ & Vecteur, \texttt{float} & nombre de particules hors de la DOF
(\textit{comptes par p�riode d'�chantillonnage})\\
TAS & Vecteur, \texttt{float} & Vitesse propre (m/s)\\
$\tau{}_1$ & Coeff., \texttt{float} & d�lai long (second)\\
$\tau{}_2$ & Coeff., \texttt{float} & d�lai court (second)\\
\end{tabarg}

\medskip
%%%%%SORTIES%%%%
\begin{tabarg}
\multicolumn{3}{l}{\titretab {Sorties}}\\
\hline
$A_c$ & Vecteur, \texttt{float} & Activit� (proportion : $s$/$s$)\\
\end{tabarg}


%%%%%FORMULE%%%%
\itemstyle{Formule}

%% Formule de l'operateur : fssp_activ_calc
%% Dans ce fichier, les �ventuels accents sont cod�s
%% en ISO-8859-15 (aka latin9)
%% La ligne suivante donne trois exemples 
%% e accent aigu : �, e accent grave : �, e-dans-l'o : �
\[
\tau_T =  \frac{\Phi_{\text{Beam}} }{ \text{TAS} }
\]
\[
A_c = Strobe.(\tau_1 + \tau_T) + Fast\_Reset . (\tau_2 +\tau_T)
\]



%%%%%COMMENTAIRE%%%%
\itemstyle{Commentaire~:}
La formule de la note 24 du NCAR ne tient pas compte du temps
de transit moyen des gouttes dans le faisceau ($\tau_T$). Les
d�lais longs et courts ($\tau_1$ et $\tau_2$) sont des d�lais 
de traitement �lectronique, il est  donc logique de leur rajouter
le temps de transit moyen.

%%%%%AUTEUR%%%%
\itemstyle{Auteur~:}
CNRM/GMEI/TRAMM

%%%%%REFERENCE%%%%
\itemstyle{R\'ef\'erences~:}
Darrel Baumgardner, Bulletin NCAR/RAF n�24,
\textit{Airborne Measurements for Cloud Microphysics}.
\cite{RAF-24}
